\chapter*{Abstract}

Green AI research has established that the energy consumption and emissions of machine learning are measurable, but its empirical scope has remained tied to large-scale deep learning on text and image data. The classical models that dominate deployed tabular classification lack a comparable baseline. This thesis addresses that gap through a systematic benchmark of five classifiers, Logistic Regression, Random Forest, XGBoost on \acrshort{cpu} and \acrshort{gpu}, and a neural network, on three tabular datasets spanning four orders of magnitude in size, from 178 to 11 million instances. The scope extends beyond the single training run to cover hyperparameter search, temporal variation in electricity grid carbon intensity and the inference break-even point at which cumulative deployment cost overtakes the training carbon footprint. To address the systematic undercounting of software-based estimators, all \acrshort{cpu} energy figures are derived from direct hardware-sensor readings via \texttt{CodeCarbon}. A composite metric, \acrlong{cf1}, is introduced to express \acrlong{co2eq} emissions per unit of predictive performance.

The results show that ecological efficiency is not a property of any single model but emerges from the fit between its computational structure, the hardware it runs on and the volume of data it processes. No model leads across all three datasets. \acrshort{gpu} acceleration becomes the more efficient choice for XGBoost only above approximately 80{,}000 instances. Below that threshold device initialization overhead dominates and emissions exceed those of the \acrshort{cpu} variant. A constrained \texttt{Optuna} search for the \acrshort{mlp} on HIGGS emitted roughly 54 times the carbon emissions of the final training run it enabled. Further analysis reveals that direct hardware measurement exceeded raw \texttt{CodeCarbon} figures by a median factor of 1.92, and shifting training away from winter evenings towards summer midday reduced absolute emissions by up to 70 percent without altering a single line of code. The contribution of this work is not a ranked list of preferred classifiers but a corrected measurement methodology, a context-sensitive efficiency metric and a lifecycle accounting framework that together make ecological model selection answerable for a given deployment context rather than as a general question.
